\documentclass[11pt,a4paper]{article}
\usepackage[T1]{fontenc}
\usepackage{lmodern}
\usepackage[margin=27mm]{geometry}
\usepackage{amsmath,amssymb,amsthm,mathtools}
\usepackage{microtype}
\usepackage{booktabs,array}
\usepackage{enumitem}
\usepackage[hidelinks]{hyperref}
\setlength{\parskip}{5pt plus 1pt}
\setlength{\parindent}{0pt}
\setlength{\emergencystretch}{3em}
\setlist{itemsep=2pt,topsep=3pt}
\raggedbottom
\newtheorem{theorem}{Theorem}[section]
\newtheorem{lemma}[theorem]{Lemma}
\newtheorem{corollary}[theorem]{Corollary}
\newtheorem{proposition}[theorem]{Proposition}
\theoremstyle{definition}
\newtheorem{definition}[theorem]{Definition}
\newtheorem{example}[theorem]{Example}
\theoremstyle{remark}
\newtheorem{remark}[theorem]{Remark}
\newcommand{\Z}{\mathbb{Z}}
\newcommand{\N}{\mathbb{N}}
\newcommand{\Q}{\mathbb{Q}}
\newcommand{\F}{\mathbb{F}}
\newcommand{\ps}[1]{#1[[x]]}
\newcommand{\coeff}[2]{[x^{#1}]\,#2}
\newcommand{\hl}[1]{\textbf{#1}}
\newcommand{\para}[1]{\par\medskip\noindent\textbf{#1}\enspace}
\hypersetup{pdftitle={Remainders modulo 4 for the series with A(x/A(x)) = theta3(x)},pdfauthor={Dev D. Goyal},pdfsubject={OEIS A378580}}
\title{Remainders modulo $4$ for the series with $A\bigl(x/A(x)\bigr)=\theta_3(x)$\\[6pt]\large A proof of a conjecture in OEIS A378580}
\author{Dev D. Goyal}
\date{October 2026}
\begin{document}
\maketitle
\begin{abstract}
Let $\theta(x)=1+2x+2x^4+2x^9+2x^{16}+\cdots$ be the Jacobi theta series $\theta_3$, whose coefficient of $x^n$ is $2$ when $n$ is a
positive perfect square and $0$ otherwise. Entry A378580 of the On-Line Encyclopedia of Integer Sequences (OEIS),
contributed by Paul~D.~Hanna, lists the coefficients $a(n)$ of the power series $A(x)=1+2x+4x^2+16x^3+98x^4+\cdots$
determined by $A\bigl(x/A(x)\bigr)=\theta(x)$. The entry conjectures that for $n>0$ the coefficient $a(n)$ leaves remainder $2$ on
division by $4$ when $n$ is a perfect square, and is divisible by $4$ otherwise. We prove this conjecture. In fact,
modulo $4$ the whole series $A(x)$ agrees with $\theta(x)$, and the same is true for the companion equation
$A\bigl(x\,A(x)\bigr)=\theta(x)$ of entry A378581. The proof rests on one observation: for every series $H$, the
reduction of $2H(P)$ modulo $4$ depends only on the reduction of $P$ modulo $2$. The results were checked by computer and proved
formally in the Lean proof assistant.
\end{abstract}
\begin{center}\begin{minipage}{0.88\linewidth}\small
\textbf{How this paper was produced.} This paper reports the result of an AI research effort directed by
Dev~D.~Goyal. ChatGPT (OpenAI), using the GPT-6 model with Pro reasoning and working autonomously from prompts by
Dev~D.~Goyal, found the proof, wrote its Lean~4 formalization, and wrote a first draft of the paper.
Dev~D.~Goyal chose the problems and asked for proofs, formal verification and revisions, but did not supply proof
ideas. Claude Opus~5.5, used with Claude Code, then checked the work independently: it rebuilt the formal proofs from
source and confirmed that they rely only on Lean's three standard axioms, recomputed the sequence from its definition
compared it with the published terms and tested every claim numerically, and read every step of the proofs. It also wrote the present
version of this paper.
Dev~D.~Goyal has reviewed the proofs and is responsible for them; this review is not an independent expert validation.
The paper has not been peer reviewed.
The section ``About this paper'' at the end gives the details.
\end{minipage}\end{center}

\section{Introduction}

\subsection{The sequence}\label{sec:sequence}
The \emph{On-Line Encyclopedia of Integer Sequences} (OEIS)~\cite{oeisdb} is a large public database of integer
sequences. Many
entries also record \emph{conjectures}: patterns that the contributor observed in the data but did not prove.

This paper is about entry A378580~\cite{oeis}, contributed by Paul~D.~Hanna in January 2025. It starts from the
\emph{theta series}
\[
\theta(x)=1+2\sum_{j\ge1}x^{j^2}=1+2x+2x^4+2x^9+2x^{16}+2x^{25}+\cdots .
\]
Another way to write it is $\theta(x)=\sum_{k}x^{k^2}$, where $k$ runs over all integers, positive, negative and zero:
each positive square $j^2$ appears twice, once from $k=j$ and once from $k=-j$. So the coefficient of $x^n$ in
$\theta(x)$ is the number of integers whose square is $n$. This series is one of the classical theta functions studied
by Jacobi, usually denoted $\theta_3$; the OEIS writes it as \texttt{theta\_3}, and we write $\theta$.

The entry defines a power series $A(x)=a(0)+a(1)x+a(2)x^2+\cdots$, with $a(0)=1$, by the equation
\begin{equation}\label{eq:entry}
A\bigl(x/A(x)\bigr)=\theta(x) .
\end{equation}
Here $x/A(x)$ means $x$ times the multiplicative inverse $1/A(x)$, which exists because $a(0)=1$
(Section~\ref{sec:background}). It is a series with constant term $0$, and the left side of \eqref{eq:entry} substitutes
it for $x$ in $A$. The entry reads:
\begin{quote}\small\ttfamily\raggedright G.f. A(x) satisfies: A(x/A(x)) = theta\_3(x) = 1 + 2*Sum\_\{n>=1\} x\^{}(n\^{}2).\par\smallskip Conjecture: for n > 0, a(n) == 2 (mod 4) iff n is square, else a(n) is divisible by 4 if n is nonsquare.\par\smallskip What is the radius of convergence of g.f. A(x)?\par\smallskip\hfill\textrm{--- OEIS A378580, Paul D. Hanna, January 2025}\end{quote}

Here ``G.f.''\ stands for \emph{generating function}, the series whose coefficients are the terms of the sequence.

Equation~\eqref{eq:entry} determines the coefficients one at a time. Write $A=1+a(1)x+a(2)x^2+\cdots$. Then
$1/A=1-a(1)x+\cdots$, so $x/A=x-a(1)x^2+\cdots$, and
\[
A(x/A)=1+a(1)\bigl(x-a(1)x^2\bigr)+a(2)x^2+\cdots=1+a(1)\,x+\bigl(a(2)-a(1)^2\bigr)x^2+\cdots .
\]
Comparing with $\theta(x)=1+2x+0x^2+\cdots$ gives $a(1)=2$ and $a(2)=4$. In the same way, the coefficient
of $x^n$ in \eqref{eq:entry} determines $a(n)$ once the earlier coefficients are known (Lemma~\ref{lem:unique}). So
there is exactly one solution, and all its coefficients are integers. Its first terms are
1, 2, 4, 16, 98, 756, 6848, 70000, 787056, 9569826, \dots{} (starting from $n=0$).

\subsection{The question}
The conjecture says that for $n>0$, the coefficient $a(n)$ leaves remainder $2$ on division by $4$ when $n$ is a
perfect square ($n=1,4,9,16,\dots$), and remainder $0$ when $n$ is not a square. (Section~\ref{sec:background} recalls
the notation $a\equiv b\pmod 4$ for ``$a$ and $b$ leave the same remainder on division by $4$''.) The table shows the
remainders of the first terms. The indices in bold are the positive squares.
\begin{center}\small\setlength{\tabcolsep}{3.2pt}
\begin{tabular}{>{\raggedright\arraybackslash}p{6.4em}>{\centering\arraybackslash}p{1.15em}>{\centering\arraybackslash}p{1.15em}>{\centering\arraybackslash}p{1.15em}>{\centering\arraybackslash}p{1.15em}>{\centering\arraybackslash}p{1.15em}>{\centering\arraybackslash}p{1.15em}>{\centering\arraybackslash}p{1.15em}>{\centering\arraybackslash}p{1.15em}>{\centering\arraybackslash}p{1.15em}>{\centering\arraybackslash}p{1.15em}>{\centering\arraybackslash}p{1.15em}>{\centering\arraybackslash}p{1.15em}>{\centering\arraybackslash}p{1.15em}>{\centering\arraybackslash}p{1.15em}>{\centering\arraybackslash}p{1.15em}>{\centering\arraybackslash}p{1.15em}}
\toprule
$n$ & 0 & \hl{1} & 2 & 3 & \hl{4} & 5 & 6 & 7 & 8 & \hl{9} & 10 & 11 & 12 & 13 & 14 & 15 \\ \midrule
$a(n)\bmod 4$ & 1 & \hl{2} & 0 & 0 & \hl{2} & 0 & 0 & 0 & 0 & \hl{2} & 0 & 0 & 0 & 0 & 0 & 0 \\
\bottomrule
\end{tabular}
\\[4pt]
\begin{tabular}{>{\raggedright\arraybackslash}p{6.4em}>{\centering\arraybackslash}p{1.15em}>{\centering\arraybackslash}p{1.15em}>{\centering\arraybackslash}p{1.15em}>{\centering\arraybackslash}p{1.15em}>{\centering\arraybackslash}p{1.15em}>{\centering\arraybackslash}p{1.15em}>{\centering\arraybackslash}p{1.15em}>{\centering\arraybackslash}p{1.15em}>{\centering\arraybackslash}p{1.15em}>{\centering\arraybackslash}p{1.15em}>{\centering\arraybackslash}p{1.15em}>{\centering\arraybackslash}p{1.15em}>{\centering\arraybackslash}p{1.15em}>{\centering\arraybackslash}p{1.15em}>{\centering\arraybackslash}p{1.15em}>{\centering\arraybackslash}p{1.15em}}
\toprule
$n$ & \hl{16} & 17 & 18 & 19 & 20 & 21 & 22 & 23 & 24 & \hl{25} & 26 & 27 & 28 & 29 & 30 & 31 \\ \midrule
$a(n)\bmod 4$ & \hl{2} & 0 & 0 & 0 & 0 & 0 & 0 & 0 & 0 & \hl{2} & 0 & 0 & 0 & 0 & 0 & 0 \\
\bottomrule
\end{tabular}
\\[4pt]
\begin{tabular}{>{\raggedright\arraybackslash}p{6.4em}>{\centering\arraybackslash}p{1.15em}>{\centering\arraybackslash}p{1.15em}>{\centering\arraybackslash}p{1.15em}>{\centering\arraybackslash}p{1.15em}>{\centering\arraybackslash}p{1.15em}>{\centering\arraybackslash}p{1.15em}>{\centering\arraybackslash}p{1.15em}>{\centering\arraybackslash}p{1.15em}>{\centering\arraybackslash}p{1.15em}>{\centering\arraybackslash}p{1.15em}>{\centering\arraybackslash}p{1.15em}>{\centering\arraybackslash}p{1.15em}>{\centering\arraybackslash}p{1.15em}>{\centering\arraybackslash}p{1.15em}>{\centering\arraybackslash}p{1.15em}>{\centering\arraybackslash}p{1.15em}}
\toprule
$n$ & 32 & 33 & 34 & 35 & \hl{36} & 37 & 38 & 39 & 40 & 41 & 42 & 43 & 44 & 45 & 46 & 47 \\ \midrule
$a(n)\bmod 4$ & 0 & 0 & 0 & 0 & \hl{2} & 0 & 0 & 0 & 0 & 0 & 0 & 0 & 0 & 0 & 0 & 0 \\
\bottomrule
\end{tabular}
\end{center}
Every term in the table follows the conjectured pattern. The entry also asks for the radius of
convergence of $A(x)$, a question about how fast the coefficients grow; we return to it in
Section~\ref{sec:discussion}.

\subsection{What this paper proves}
\textbf{Main result} (Theorem~\ref{thm:main}). \emph{Modulo $4$, the series $A(x)$ is equal to $\theta(x)$: every
coefficient $a(n)$ leaves the same remainder on division by $4$ as the corresponding coefficient of $\theta(x)$.}

Since the coefficients of $\theta(x)$ are $1$ at $n=0$, $2$ at the positive squares and $0$ elsewhere, this is exactly
Hanna's conjecture (Corollary~\ref{cor:conj}). The theorem also covers the companion entry A378581~\cite{oeisother},
defined by $A\bigl(x\,A(x)\bigr)=\theta(x)$, which has the same conjecture. Both entries are treated by the same theorem and
the same proof; the two papers differ only in the description of the entry.

\subsection{What was known}
The entry records several formulas due to Hanna, for example $A(x)=\theta\bigl(B(x)\bigr)$, where $B$ is the
compositional inverse of $x/A(x)$. These formulas do not seem to give the remainders directly. We are
not aware of an earlier proof of the conjecture.

\subsection{Motivation}
Theta series are classical objects of number theory. For instance, the coefficient of $x^n$ in
$\theta(x)^k$ is the number of ways to write $n$ as a sum of $k$ squares of integers, and classical theorems of Jacobi
give exact formulas for these numbers when $k=2$ and $k=4$~\cite{hardywright}. Hanna's entries explore what happens when
such classical series are combined with substitution. The result here says that, modulo $4$, the substitution in
\eqref{eq:entry} has no effect: the solution is congruent to $\theta$ itself.

The argument is general (Section~\ref{sec:discussion}): it applies to every target series of the form $1+2S$, and it
treats both entries at once. Its main tool, Lemma~\ref{lem:double}, passes from information modulo $2$ to information
modulo $4$; it is also used in other papers of this collection.

\subsection{How to read this paper}
Section~\ref{sec:background} explains the few notions we need: power series, substitution, and remainders.
Section~\ref{sec:results} states the results precisely, and Section~\ref{sec:proofs} proves them.
Section~\ref{sec:checks} describes how the results were checked, by computer and by a formal proof in Lean.
Section~\ref{sec:discussion} discusses the scope of the result, a generalization, and a further congruence between the
two companion entries.
\section{Background}\label{sec:background}
This section collects the definitions and facts used later. Readers who know formal power series and congruences can skip it and return when needed. Standard references for formal power series are~\cite{niven,stanley}.
\subsection{Formal power series}
A \emph{formal power series} is an expression
\[
U(x)=u_0+u_1x+u_2x^2+u_3x^3+\cdots ,
\]
determined by its sequence of coefficients $u_0,u_1,u_2,\dots$; see~\cite{niven,stanley}. No number is ever
substituted for $x$, so questions of convergence do not arise. We write $\coeff{n}{U}=u_n$ for the coefficient of $x^n$ in $U$; two series are equal
when all their coefficients are equal.

Series are added coefficient by coefficient and multiplied like polynomials:
\[
\coeff{n}{(UV)}=u_0v_n+u_1v_{n-1}+\cdots+u_nv_0 .
\]
Each coefficient of a product is a finite sum. For example, $(1-x)(1+x+x^2+x^3+\cdots)=1$, so $1+x+x^2+\cdots$ is
the \emph{multiplicative inverse} of $1-x$, written $\frac{1}{1-x}$. More generally, every series $D$ with constant term
$1$ has a unique multiplicative inverse $\frac1D$, whose coefficients are determined one at a time by $D\cdot\frac1D=1$
using only addition, subtraction and multiplication. Whenever a quotient of power series appears in this paper, its
denominator has constant term $1$, so the quotient is again a power series.

A \emph{commutative ring} is a set with an addition, a subtraction and a multiplication that satisfy the usual rules
of arithmetic; examples are $\Z$, $\Q$ and the integers modulo $m$. For a commutative ring $R$ we write $\ps{R}$ for
the set of power series with coefficients in $R$; for instance, writing $\F_2$ for the integers modulo $2$, $\ps{\F_2}$
consists of the series with coefficients modulo $2$. We write
$U(0)$ for the constant term $u_0$ of $U$, and $U=O(x^d)$ when the coefficients of $1,x,\dots,x^{d-1}$ in $U$ are all
zero. Thus $U-V=O(x^d)$ means that $U$ and $V$ have the same coefficient of $x^n$ for every $n<d$. A series with
constant term $0$ and coefficient of $x$ equal to $1$, that is, a series of the form $x+u_2x^2+u_3x^3+\cdots$, is
called \emph{normalized}.

\subsection{Substituting one series into another}
If $V(x)$ has constant term $0$, we can substitute $V$ for $x$ in $U$:
\[
U(V(x))=u_0+u_1V(x)+u_2V(x)^2+u_3V(x)^3+\cdots .
\]
This is called \emph{composition}, also written $U\circ V$. It makes sense because $V(x)^k$ has no terms below $x^k$,
so only the finitely many terms with $k\le n$ affect the coefficient of $x^n$. For example, if $U=x+x^2$ and $V=x^2$, then
$U(V(x))=x^2+x^4$. Composition is associative: $U\circ(V\circ W)=(U\circ V)\circ W$.

Substituting a series into itself repeatedly gives its \emph{iterates}: $U^{\circ2}(x)=U(U(x))$,
$U^{\circ3}(x)=U(U(U(x)))$, and so on, with $U^{\circ0}(x)=x$ and $U^{\circ1}=U$. They must not be confused with
\emph{powers} $U^2=U\cdot U$, $U^3=U\cdot U\cdot U$. (The OEIS sometimes writes $A^k(x)$ for the $k$-th iterate. In
this paper the small circle always means an iterate, and a plain exponent always means a power.)

\subsection{Remainders and congruences}
For an integer $m\ge2$ we write $a\equiv b\pmod m$ (``$a$ is congruent to $b$ modulo $m$'') when $a-b$ is divisible by
$m$, that is, when $a$ and $b$ leave the same remainder after division by $m$. The remainder of $a$ on division by
$m$ is called its \emph{residue} modulo $m$; it is one of $0,1,\dots,m-1$.

Addition and multiplication are compatible with congruences, so the residues $0,1,\dots,m-1$, added and multiplied
modulo $m$, form a ring, written $\Z/m\Z$. When $m=p$ is a prime we also write $\F_p$. In $\F_2=\{0,1\}$ we have $1+1=0$,
so addition and subtraction coincide.

For power series with integer coefficients, $U\equiv V\pmod m$ means that \emph{every} coefficient of $U-V$ is
divisible by $m$. Equivalently, the two series become equal when every coefficient is replaced by its residue; we call
this \emph{reducing} the series modulo $m$. The key fact is the following.

\emph{Reduction modulo $m$ is compatible with all the operations above.} To reduce a sum, a product, a substitution, or
an inverse, we may reduce the ingredients first and then compute in $\Z/m\Z$. This holds because every
coefficient of the result is obtained from finitely many coefficients of the ingredients by adding, subtracting and
multiplying, and these operations are compatible with remainders.

\para{Rings.} Some statements below hold over every commutative ring, for example $\Z$, $\Q$ and $\Z/m\Z$. No
division is used in them; this matters because in $\Z/m\Z$ division is not always possible (for instance, $2$ has no inverse in
$\Z/4\Z$).

\subsection{The method}
The main proofs follow three steps.
\begin{enumerate}
\item \emph{Uniqueness.} Show that the defining equation determines the coefficients one at a time, using only
addition, subtraction and multiplication. Then the equation has exactly one solution of the required shape, in every
commutative ring, including $\Z/4\Z$.
\item \emph{Reduction.} Reducing the integer solution modulo $4$ gives a solution of the reduced equation,
because reduction is compatible with all operations.
\item \emph{A candidate solution.} Find an explicit series that satisfies the reduced equation. By uniqueness, it is
the reduction of the integer solution, and its coefficients can be read off.
\end{enumerate}
Finding the candidate in step~3 is usually the main difficulty.
\subsection{Doubling modulo 4}
\begin{lemma}\label{lem:double}
Let $U,V$ be series with coefficients in $\Z/4\Z$. Then $2U=2V$ if and only if $U$ and $V$ agree modulo $2$.
Consequently, if $P$ and $Q$ have constant term $0$ and $P\equiv Q\pmod 2$, then $2H(P)=2H(Q)$ in $\Z/4\Z$ for every
series $H$.
\end{lemma}

\begin{proof}
For residues $u,v$ modulo $4$, $2u=2v$ means that $4$ divides $2(u-v)$, that is, $2$ divides $u-v$. Apply this to
every coefficient. For the second claim, $H(P)$ and $H(Q)$ agree modulo $2$ because reduction is compatible with
substitution.
\end{proof}
In other words, \emph{twice a series, modulo $4$, depends only on the series modulo $2$.} This lets us use information modulo
$2$ to get information modulo $4$, without ever dividing by $2$ in $\Z/4\Z$, which is impossible.

\subsection{A one-term Taylor formula}
\begin{lemma}\label{lem:taylor}
Let $\varepsilon$ be an element of a commutative ring with $\varepsilon^2=0$. If $u$ and $v$ are series with constant
term $0$, then for every series $S$
\[
S(u+\varepsilon v)=S(u)+\varepsilon\,v\,S'(u),
\]
where $S'=\sum_{n\ge1}n\,s_nx^{n-1}$ is the formal derivative of $S$.
\end{lemma}

\begin{proof}
By the binomial theorem, $(u+\varepsilon v)^j=u^j+j\,\varepsilon v\,u^{j-1}$, because all terms with
$\varepsilon^2$ or higher powers vanish. Multiply by $s_j$ and add up over $j$.
\end{proof}
An example of such an $\varepsilon$: in $\Z/8\Z$, the element $4$ satisfies $4^2=16=0$.

\section{Results}\label{sec:results}

\begin{theorem}\label{thm:main}
Each of the two equations
\[
A\bigl(x/A(x)\bigr)=\theta(x)\qquad\text{and}\qquad A\bigl(x\,A(x)\bigr)=\theta(x)
\]
has exactly one solution $A$ with integer coefficients and $A(0)=1$. These solutions are the generating functions of
A378580 and A378581, respectively. Both satisfy
\[
A(x)\equiv\theta(x)=1+2\sum_{j\ge1}x^{j^2}\pmod 4 .
\]
\end{theorem}

\begin{corollary}[Hanna's conjecture]\label{cor:conj}
For the sequence $a(n)$ of OEIS A378580 and every $n>0$: if $n$ is a perfect square, then $a(n)\equiv2\pmod4$;
otherwise $a(n)\equiv0\pmod4$, that is, $4$ divides $a(n)$.
\end{corollary}

\begin{proof}
Compare the coefficients of $x^n$ on both sides of the congruence in Theorem~\ref{thm:main}. For $n>0$ the coefficient
of $x^n$ in $\theta(x)$ is $2$ if $n$ is a perfect square and $0$ otherwise.
\end{proof}

\section{Proofs}\label{sec:proofs}

\para{The idea in brief.} Work modulo $4$ and write $\theta=1+2S$, with $S=x+x^4+x^9+\cdots$. Then
$\theta^2=1+4S+4S^2=1$, so $\theta$ is its own multiplicative inverse, and $x/\theta=x\theta$. Both equations therefore
ask whether $\theta(x\theta)=\theta$. Now $x\theta=x+2xS$ differs from $x$ by a multiple of $2$. Substituting it into
$\theta=1+2S$ changes $S$ by a multiple of $2$, so it changes $2S$ by a multiple of $4$, that is, not at all modulo $4$.
So $\theta$ solves both equations modulo $4$, and by uniqueness it is the reduction of both solutions.

\begin{lemma}[Existence and uniqueness]\label{lem:unique}
Let $R$ be a commutative ring and let $T$ be a series with coefficients in $R$ and constant term $1$. Each of the
equations
\[
U\bigl(x\,U(x)\bigr)=T(x)\qquad\text{and}\qquad U\bigl(x/U(x)\bigr)=T(x)
\]
has exactly one solution $U$ with coefficients in $R$ and constant term $1$.
\end{lemma}

\begin{proof}
Write $U=1+u_1x+u_2x^2+\cdots$ and $T=1+t_1x+t_2x^2+\cdots$.

\emph{The product form.} Since $U(y)=\sum_ku_ky^k$, we have $U(xU)=\sum_{k\ge0}u_k\,x^kU^k$, and the term with index $k$
has no terms below degree $k$. So for $n\ge1$ the coefficient of $x^n$ in $U(xU)$ is
\[
u_n\cdot\coeff{0}{U^n}+\sum_{k=1}^{n-1}u_k\,\coeff{n-k}{U^k}
=u_n+\bigl(\text{an expression in }u_1,\dots,u_{n-1}\bigr),
\]
because $\coeff{0}{U^n}=1$, and $\coeff{n-k}{U^k}$ involves only $u_1,\dots,u_{n-k}$, where $n-k\le n-1$. (The term
$k=0$ is the constant $1$.) Setting this coefficient equal to $t_n$ determines $u_n$ from $u_1,\dots,u_{n-1}$, using only
addition, subtraction and multiplication. So the coefficients are determined one at a time: there is exactly one
solution.

\emph{The quotient form.} Put $W=1/U$, the multiplicative inverse of $U$, so that $x/U=xW$. Its coefficients are
determined by $UW=1$, that is, $w_j=-(u_1w_{j-1}+u_2w_{j-2}+\cdots+u_jw_0)$, so $w_j$ depends only on $u_1,\dots,u_j$.
Now $U(xW)=\sum_ku_k\,x^kW^k$, and exactly as before, the coefficient of $x^n$ is $u_n$ plus an expression in
$w_1,\dots,w_{n-1}$ and $u_1,\dots,u_{n-1}$, hence in $u_1,\dots,u_{n-1}$ alone. So again the coefficients are determined
one at a time.
\end{proof}

No division was used in this proof. So, as explained in Section~\ref{sec:background}, reducing the integer solution
modulo $4$ gives the solution over $\Z/4\Z$, and to find that reduction it is enough to find \emph{some} solution over
$\Z/4\Z$.

\begin{proof}[Proof of Theorem~\ref{thm:main}]
The first two claims are Lemma~\ref{lem:unique} with $R=\Z$ and $T=\theta$. For the congruence, work over $\Z/4\Z$ and
write $\theta=1+2S$ with $S=\sum_{j\ge1}x^{j^2}$.

\emph{Step 1.} $\theta^2=1+4S+4S^2=1$, because $4=0$ in $\Z/4\Z$. So $1/\theta=\theta$, and therefore
$x/\theta=x\theta$.

\emph{Step 2.} The series $x\theta=x+2xS$ agrees with $x$ modulo $2$, and both have constant term $0$. By
Lemma~\ref{lem:double}, $2S(x\theta)=2S(x)$. Hence
\[
\theta(x\theta)=1+2S(x\theta)=1+2S(x)=\theta(x).
\]

\emph{Step 3.} By Steps 1 and 2, $\theta(x/\theta)=\theta(x\theta)=\theta$. So $U=\theta$ solves both equations of
Lemma~\ref{lem:unique} over $\Z/4\Z$, with $T=\theta$. By uniqueness, the reduction modulo $4$ of each integer solution
is $\theta$.
\end{proof}

\section{Checking the results}\label{sec:checks}

A mathematical proof can contain mistakes, so the results were checked in two independent ways.

\para{Numerical check.}
We computed $A$ from the defining equation~\eqref{eq:entry}, with exact integer arithmetic, through degree
$199$. This reproduces all 24 terms printed in the entry, and for every $0<n<200$ the remainder of
$a(n)$ modulo $4$ is $2$ when $n$ is a perfect square and $0$ otherwise. The companion sequence A378581 was
checked in the same way.
The computation was done with the PARI/GP computer algebra system, using code that was written separately from the
formal proofs (file \texttt{verification/independent\_checks.gp} in the accompanying repository~\cite{repo}).
A computation can only test finitely many cases, so it is a check on the definitions and on the indexing, not a proof.

\para{Formal proof in Lean.}
Lean~\cite{lean} is a \emph{proof assistant}: a program in which mathematical statements and their proofs are written
in a precise formal language, and every step is checked mechanically down to the basic rules of logic. The theorems of Section~\ref{sec:results} are proved in Lean~4 (version 4.35.0) on top of
Mathlib~\cite{mathlib}, Lean's library of formalized mathematics. The statements are collected in the file \texttt{lean/Entries/A378580.lean} of the repository,
in the namespace \texttt{OEIS.A378580}. In the Lean code, $\theta$ is called \texttt{theta}. The prefix \texttt{minus} in the
names of the declarations refers to the quotient equation $A(x/A)=\theta$; the corresponding declarations for A378581 carry the prefix
\texttt{plus}. The quotient $x/A$ is written \texttt{X * invOfUnit A 1}, the
product of $x$ with the multiplicative inverse of $A$. The general constructions, for an arbitrary series $T$ with constant term $1$
over any commutative ring, are in \texttt{lean/Families/ThetaSubstitution.lean}. The Lean proof of Lemma~\ref{lem:unique}
takes a slightly different route from the one above: it treats the product form through the equivalent equation
$G(G(x))=G(x)\,T(x)$ for $G=xU$, and reduces the quotient form to the product form: if $U(x/U)=T$, then $V=1/U$ satisfies $V(xV)=1/T$.
The table lists where each result of this paper is proved.
\begin{center}\small\renewcommand{\arraystretch}{1.25}
\begin{tabular}{>{\raggedright\arraybackslash}p{0.52\linewidth}>{\raggedright\arraybackslash}p{0.4\linewidth}}
\toprule
Statement in this paper & Lean theorem(s)\\
\midrule
$\theta(0)=1$, and for $n>0$ the coefficient of $x^n$ in $\theta$ is $2$ or $0$ as $n$ is or is not a square & \texttt{theta\_\allowbreak{}constant\_\allowbreak{}term}, \texttt{theta\_\allowbreak{}coefficient\_\allowbreak{}formula} \\
$A(0)=1$, and $A\cdot\bigl(x\cdot(1/A)\bigr)=x$, so the quotient is literally $x/A$ & \texttt{minus\_\allowbreak{}constant}, \texttt{quotient\_\allowbreak{}identity} \\
$A(x/A)=\theta$, and Lemma~\ref{lem:unique} for this equation & \texttt{minus\_\allowbreak{}source}, \texttt{minus\_\allowbreak{}exists\_\allowbreak{}unique}, \texttt{minus\_\allowbreak{}unique} \\
Theorem~\ref{thm:main} ($A\equiv\theta\pmod 4$) and Corollary~\ref{cor:conj} & \texttt{minus\_\allowbreak{}series\_\allowbreak{}mod\_\allowbreak{}four}, \texttt{minus\_\allowbreak{}congruence}, \texttt{minus\_\allowbreak{}square\_\allowbreak{}iff}, \texttt{minus\_\allowbreak{}four\_\allowbreak{}dvd\_\allowbreak{}iff} \\
\bottomrule
\end{tabular}
\end{center}
The Lean project was rebuilt from source, and the axioms on which each of these theorems depends were listed with
Lean's \texttt{\#print axioms} command. Every one depends only on the three standard axioms of Lean's logic
(\texttt{propext}, \texttt{Classical.choice}, \texttt{Quot.sound}); no step is left unproved. A formal proof certifies
exactly the statement that Lean checks. Whether that statement says what the OEIS entry means is a separate question,
so the Lean definitions were also compared with the definition in the OEIS entry
(Section~\ref{sec:sequence}). The principal theorems are stated once more, using only Mathlib's definitions, in the
file \texttt{lean/Challenges/Theta.lean}, and Lean's Comparator tool confirms that the
proved theorems have exactly these statements.

\section{Discussion}\label{sec:discussion}

\para{Scope.} Theorem~\ref{thm:main} describes the remainder modulo $4$ of every coefficient, so it settles the
conjecture completely, for this entry and for its companion A378581. It says nothing about the size of the
coefficients. In particular, it does not answer the entry's question about the radius of convergence of
$A(x)$, which is the reciprocal of the exponential growth rate of $|a(n)|$. That is a question about the size of
the numbers $a(n)$, which remainders cannot detect.

\para{A generalization.} The proof used only that $\theta=1+2S$ for some series $S$ with integer coefficients and
$S(0)=0$. So the same three steps show: \emph{for every such target series $T=1+2S$, the solutions of $U(xU)=T$ and of
$U(x/U)=T$ with constant term $1$ are both congruent to $T$ modulo $4$.} As a check, we confirmed this by computer for
$20$ randomly chosen targets $T$, up to degree $59$. This general form is not part of the Lean formalization.

\para{The two entries agree modulo $8$.} Neither solution agrees with $\theta$ modulo $8$: for both entries
$a(2)=\pm4$, while the coefficient of $x^2$ in $\theta$ is $0$. But the two sequences A378580 and A378581 agree with
\emph{each other} modulo $8$. More generally, for every series $T$ with integer coefficients and $T\equiv1\pmod 2$, the
solutions of $U(x/U)=T$ and $U(xU)=T$ are congruent modulo $8$. Here is the argument, which is not part of the Lean
formalization. Work over $\Z/8\Z$, and let $U$ solve $U(xU)=T$. Modulo $2$ we have $T=1$, and then $U=1$ is the
solution, so $U=1+2W$ for some series $W$. Then $U^2=1+4(W+W^2)$, and since $(1+4X)^2=1+8X+16X^2=1$ in $\Z/8\Z$, we get
$1/U^2=1+4(W+W^2)$ and
\[
\frac xU=x\,U\cdot\frac1{U^2}=x\,U+4E,\qquad E=x\,U\,(W+W^2).
\]
By Lemma~\ref{lem:taylor} with $\varepsilon=4$ (note $4^2=16=0$ in $\Z/8\Z$),
$U(xU+4E)=U(xU)+4E\,U'(xU)$. The derivative $U'=2W'$ is a multiple of $2$, so the last term is a multiple of $8$, which
is $0$. Hence $U(x/U)=U(xU)=T$: the same $U$ solves the quotient equation, and by Lemma~\ref{lem:unique} the two
solutions agree modulo $8$. The two sequences do not agree modulo $16$: already $a(2)=4$ for A378580 and $a(2)=-4$ for
A378581, which differ by $8$. A computer check confirms the agreement modulo $8$ for every $n<160$.

\para{Context.} Section~\ref{sec:background} explained the general method: reduce modulo a small number, find an explicit
solution of the reduced equation, and conclude by uniqueness. Here that solution is the target itself. The same
doubling argument (Lemma~\ref{lem:double}) appears in other papers of this collection, for example for the theta-type
sums of entries A380060 and A380680.

\section*{About this paper}\label{sec:about}
\para{How the work was done.}
This paper is one of 28 papers on conjectures from the OEIS that came out of an AI research effort directed by
Dev~D.~Goyal~\cite{repo}. The roles were as follows.
\begin{itemize}
\item \emph{Dev~D.~Goyal} started the project, chose which problems to pursue, asked for proofs, formal verification and
revisions, and decided what to publish. Dev~D.~Goyal did not supply proof strategies or intermediate arguments.
\item \emph{ChatGPT} (OpenAI), using the GPT-6 model with Pro reasoning, worked autonomously from those prompts. It
searched the OEIS for suitable conjectures, found the proofs, wrote the Lean~4 formalization, and wrote a first draft
of this paper.
\item \emph{Claude Opus~5.5}, used with Claude Code, checked the work independently. It rebuilt the Lean project from source and confirmed that the theorems listed in Section~\ref{sec:checks} rely only on
Lean's standard axioms; compared the Lean definitions with the OEIS definition; recomputed the sequence with separately
written PARI/GP code, compared it with the terms published in the OEIS, and tested every claim of the paper on many
terms; and read every step of the proofs. It then wrote the present version of the paper, expanding the
explanations and, in places, writing out an argument in more detail or by a more elementary route. The theorems of
Section~\ref{sec:results} and the proof strategies are ChatGPT's. The Discussion section, including the short arguments
in it, was written by Claude Opus~5.5; its observations were checked by computer and are not part of the Lean
formalization.
\end{itemize}
Dev~D.~Goyal has reviewed the proofs and is responsible for this paper; this review is not an independent expert
validation.

\para{Review status.} This paper has not been peer reviewed, and it has not been submitted to a journal or to the OEIS.

\para{Priority.} An OEIS entry that labels a statement as a conjecture shows only that its contributor did not know
a proof; it does not show that no proof existed. During the research only limited searches for earlier work were
made. No claim of priority is made.

\para{Corrections.} Please report any error in the issue tracker of the repository~\cite{repo}, where the Lean sources and the numerical checks are also
available.

\section*{Acknowledgements}
The sequence, its companion A378581, and the conjecture are due to Paul~D.~Hanna~\cite{oeis,oeisother}.

\begin{thebibliography}{99}\small
\raggedright
\bibitem{oeis} P.~D.~Hanna, entry A378580 (G.f.\ $A(x)$ satisfies $A\bigl(x/A(x)\bigr)=\theta(x)$), \emph{The On-Line Encyclopedia of Integer Sequences}, January 2025. \url{https://oeis.org/A378580}. Text as of October 2026.
\bibitem{oeisother} P.~D.~Hanna, entry A378581 (G.f.\ $A(x)$ satisfies $A(xA(x))=\theta_3(x)$), \emph{The On-Line Encyclopedia of Integer Sequences}, January 2025. \url{https://oeis.org/A378581}. Text as of October 2026.
\bibitem{hardywright} G.~H.~Hardy and E.~M.~Wright, \emph{An Introduction to the Theory of Numbers}, sixth edition, Oxford University Press, 2008. (Chapter~20 treats sums of two and four squares.)
\bibitem{oeisdb} N.~J.~A.~Sloane and The OEIS Foundation Inc., \emph{The On-Line Encyclopedia of Integer Sequences}. \url{https://oeis.org}.
\bibitem{stanley} R.~P.~Stanley, \emph{Enumerative Combinatorics}, Volume~2, Cambridge University Press, 1999. (Chapter~5 treats formal power series, composition and compositional inverses.)
\bibitem{niven} I.~Niven, Formal power series, \emph{American Mathematical Monthly} 76 (1969), 871--889.
\bibitem{lean} L.~de~Moura and S.~Ullrich, The Lean~4 theorem prover and programming language, in \emph{Automated Deduction -- CADE 28}, Lecture Notes in Computer Science 12699, Springer, 2021, 625--635.
\bibitem{mathlib} The mathlib Community, The Lean mathematical library, in \emph{Proceedings of the 9th ACM SIGPLAN International Conference on Certified Programs and Proofs (CPP 2020)}, 367--381. Version used: commit \texttt{56a005855ea1}.
\bibitem{repo} D.~D.~Goyal, \emph{Proofs and corrections of OEIS conjectures}, repository, 2026. \url{https://github.com/dvgyl/oeis-proofs}.
\end{thebibliography}
\end{document}
